% ============================================================
% I. INTRODUCCIÓN
% ============================================================

\section{Introducción}

El presente informe reúne los diez ejercicios prácticos de la primera
ronda del curso de Señales y Sistemas, agrupados originalmente en las
carpetas \texttt{WORK\_ONE} a \texttt{WORK\_TEN} del repositorio. Cada
carpeta implementa, de forma independiente y con su propia suite de
pruebas automatizadas, uno o más de los ejercicios propuestos en el
enunciado de la tarea: análisis espectral de audio, series de Fourier,
convolución discreta, señales de modulación de amplitud y aliasing,
transformada wavelet, transformada de Hilbert, filtros digitales de
promedio, y diseño manual de filtros IIR mediante la transformada
bilineal, entre otros.

En lugar de un único desarrollo monolítico, se optó por mantener la
estructura modular ya existente en el repositorio (un paquete Python por
tarea, con su propio \texttt{main.py}, pruebas y \texttt{readme.md}) y
consolidar aquí únicamente los resultados: por cada carpeta se incluye
una subsección con el objetivo del ejercicio, la metodología empleada, y
las figuras y valores numéricos obtenidos al ejecutar el código
correspondiente. Las figuras mostradas en este documento son, en su
totalidad, salidas reales de los scripts de cada carpeta (o, en el caso
puntual de \texttt{WORK\_ONE}, un análisis FFT ad-hoc pero igualmente
real sobre audio generado por el propio proyecto, dado que esa carpeta no
incluye un módulo de graficado propio).

La Sección~\ref{sec:resultados} presenta los resultados de cada una de
las diez carpetas, en el mismo orden en que fueron desarrolladas. La
Sección~\ref{sec:conclusiones} sintetiza las observaciones transversales
más relevantes entre ejercicios.

El código fuente completo de los diez ejercicios (implementación, pruebas
unitarias y scripts que generan cada una de las figuras de este informe)
está disponible públicamente en
\url{https://github.com/DanielAraqueStudios/Signals-frecuency}, bajo
\texttt{THEORY/FIRST\_ROUND/WORK/}.
